% GENERATED FILE. Edit canonical lessons, then run npm run generate:books.
% canonical-source-hash: fnv1a64:d118f45035647588
% canonical-lessons: BN-C77-mittha, BN-C77-holud, BN-C77-badami, BN-C77-sundor, BN-C77-din

\chapter{False, Yellow, Brown, Beautiful, Day}
\label{ch:bn-false-yellow-brown-beautiful-day}
\hlchaptermodality{77}

\begin{chapteropening}
\textbf{By the end of this chapter:} \emph{I can say false, yellow, brown, beautiful and a day.}

Complete the last lesson of chapter 77: I can say false, yellow, brown, beautiful and a day.
\end{chapteropening}

\section[\texorpdfstring{\bn{মিথ্যা} (mitthā)}{মিথ্যা (mitthā)}]{\bn{মিথ্যা} (mitthā) — false}

\label{lesson:BN-C77-mittha}



\begin{quote}
Before the new one: say the Bengali for heavy, then the Bengali for true.
\end{quote}

\subsection*{You'll want to know: \bn{মিথ্যা}}
\textbf{\bn{মিথ্যা}} — \emph{mitthā} — \textquotedblleft{}false\textquotedblright{}.

\begin{cousinweb}
\textbf{\bn{মিথ্যা}} (\emph{mitthā}) is Sanskrit \emph{mithyā}, \textquotedblleft{}wrongly, falsely\textquotedblright{}, taken over whole as a learned loan, a tatsama (CDIAL 10130). The word inherited from it is \emph{michā}, \textquotedblleft{}untrue\textquotedblright{}, which CDIAL lists.
\end{cousinweb}

\begin{grammarlens}[title={What you've built}]
One more word for this chapter.
\end{grammarlens}

\subsection*{Your turn}
\begin{itemize}
\raggedright
  \item \emph{Say it:} \emph{mitthā}
  \item \emph{Say it:} \emph{mitthā}, once more
  \item \emph{Say it:} say \emph{mitthā}
  \item \emph{Recall:} say the Bengali for heavy, then the Bengali for true, then say \emph{mitthā} again
\end{itemize}

\begin{tcolorbox}[breakable,colback=teal!4,colframe=teal!35!black,title={Before you move on}]
What does \bn{মিথ্যা} mean? (\textquotedblleft{}False\textquotedblright{}.) Say it once more.
\end{tcolorbox}
\section[\texorpdfstring{\bn{হলুদ} (hôlud)}{হলুদ (hôlud)}]{\bn{হলুদ} (hôlud) — yellow}

\label{lesson:BN-C77-holud}



\begin{quote}
Before the new one: say the Bengali for true, then the Bengali for false.
\end{quote}

\subsection*{You'll want to know: \bn{হলুদ}}
\textbf{\bn{হলুদ}} — \emph{hôlud} — \textquotedblleft{}yellow\textquotedblright{}.

\begin{cousinweb}
\textbf{\bn{হলুদ}} (\emph{hôlud}) is a tadbhava, a word inherited from Sanskrit and reshaped by sound change. It comes from \emph{haridrā}, \textquotedblleft{}turmeric\textquotedblright{} (CDIAL 13992). Bengali names the colour after the spice.
\end{cousinweb}

\begin{grammarlens}[title={What you've built}]
One more word for this chapter.
\end{grammarlens}

\subsection*{Your turn}
\begin{itemize}
\raggedright
  \item \emph{Say it:} \emph{hôlud}
  \item \emph{Say it:} \emph{hôlud}, once more
  \item \emph{Say it:} say \emph{hôlud}
  \item \emph{Recall:} say the Bengali for true, then the Bengali for false, then say \emph{hôlud} again
\end{itemize}

\begin{tcolorbox}[breakable,colback=teal!4,colframe=teal!35!black,title={Before you move on}]
What does \bn{হলুদ} mean? (\textquotedblleft{}Yellow\textquotedblright{}.) Say it once more.
\end{tcolorbox}
\section[\texorpdfstring{\bn{বাদামি} (bādāmi)}{বাদামি (bādāmi)}]{\bn{বাদামি} (bādāmi) — brown}

\label{lesson:BN-C77-badami}



\begin{quote}
Before the new one: say the Bengali for false, then the Bengali for yellow.
\end{quote}

\subsection*{You'll want to know: \bn{বাদামি}}
\textbf{\bn{বাদামি}} — \emph{bādāmi} — \textquotedblleft{}brown\textquotedblright{}.

\begin{cousinweb}
\textbf{\bn{বাদামি}} (\emph{bādāmi}) is \emph{bādām}, \textquotedblleft{}almond\textquotedblright{}, with the ending \emph{-i}: \textquotedblleft{}almond-coloured\textquotedblright{}. \emph{Bādām} is usually traced to Persian \emph{bādām}, \textquotedblleft{}almond\textquotedblright{}.
\end{cousinweb}

\begin{grammarlens}[title={What you've built}]
One more word for this chapter.
\end{grammarlens}

\subsection*{Your turn}
\begin{itemize}
\raggedright
  \item \emph{Say it:} \emph{bādāmi}
  \item \emph{Say it:} \emph{bādāmi}, once more
  \item \emph{Say it:} say \emph{bādāmi}
  \item \emph{Recall:} say the Bengali for false, then the Bengali for yellow, then say \emph{bādāmi} again
\end{itemize}

\begin{tcolorbox}[breakable,colback=teal!4,colframe=teal!35!black,title={Before you move on}]
What does \bn{বাদামি} mean? (\textquotedblleft{}Brown\textquotedblright{}.) Say it once more.
\end{tcolorbox}
\section[\texorpdfstring{\bn{সুন্দর} (sundôr)}{সুন্দর (sundôr)}]{\bn{সুন্দর} (sundôr) — beautiful}

\label{lesson:BN-C77-sundor}



\begin{quote}
Before the new one: say the Bengali for yellow, then the Bengali for brown.
\end{quote}

\subsection*{You'll want to know: \bn{সুন্দর}}
\textbf{\bn{সুন্দর}} — \emph{sundôr} — \textquotedblleft{}beautiful\textquotedblright{}.

\begin{cousinweb}
\textbf{\bn{সুন্দর}} (\emph{sundôr}) is Sanskrit \emph{sundara}, \textquotedblleft{}beautiful\textquotedblright{}, taken over whole as a learned loan, a tatsama (CDIAL 13474). The inherited form is \emph{sũdar}, which CDIAL lists; CDIAL asks whether \emph{sundara} is itself a Middle Indo-Aryan form of \emph{sūnara}, \textquotedblleft{}excellent\textquotedblright{}.
\end{cousinweb}

\begin{grammarlens}[title={What you've built}]
One more word for this chapter.
\end{grammarlens}

\subsection*{Your turn}
\begin{itemize}
\raggedright
  \item \emph{Say it:} \emph{sundôr}
  \item \emph{Say it:} \emph{sundôr}, once more
  \item \emph{Say it:} say \emph{sundôr}
  \item \emph{Recall:} say the Bengali for yellow, then the Bengali for brown, then say \emph{sundôr} again
\end{itemize}

\begin{tcolorbox}[breakable,colback=teal!4,colframe=teal!35!black,title={Before you move on}]
What does \bn{সুন্দর} mean? (\textquotedblleft{}Beautiful\textquotedblright{}.) Say it once more.
\end{tcolorbox}
\section[\texorpdfstring{\bn{দিন} (din)}{দিন (din)}]{\bn{দিন} (din) — a day}

\label{lesson:BN-C77-din}



\begin{quote}
Before the new one: say the Bengali for brown, then the Bengali for beautiful.
\end{quote}

\subsection*{You'll want to know: \bn{দিন}}
\textbf{\bn{দিন}} — \emph{din} — \textquotedblleft{}a day\textquotedblright{}.

\begin{cousinweb}
\textbf{\bn{দিন}} (\emph{din}) is inherited from Sanskrit \emph{dina}, \textquotedblleft{}day\textquotedblright{}, and has come down with almost no change (CDIAL 6328).
\end{cousinweb}

\begin{grammarlens}[title={What you've built}]
One more word for this chapter.
\end{grammarlens}

\subsection*{Your turn}
\begin{itemize}
\raggedright
  \item \emph{Say it:} \emph{din}
  \item \emph{Say it:} \emph{din}, once more
  \item \emph{Say it:} say \emph{din}
  \item \emph{Recall:} say the Bengali for brown, then the Bengali for beautiful, then say \emph{din} again
\end{itemize}

\begin{tcolorbox}[breakable,colback=teal!4,colframe=teal!35!black,title={Before you move on}]
What does \bn{দিন} mean? (\textquotedblleft{}A day\textquotedblright{}.) Say it once more.
\end{tcolorbox}
